\section{幻觉与 jailbreak：多个回路如何竞争}

并行计算不只带来能力，也带来冲突。模型可以同时识别实体、准备答案、激活“不知道”或 refusal 倾向；最终输出取决于哪些路径被抑制、哪些路径更强。本章把 hallucination 与 jailbreak 解释为局部竞争机制，而不是一个全局开关。

从加法转到安全行为时，必须先改变读图标准：算术图的目标 token 有明确真值，而“应该回答、拒绝还是表达不确定”还受知识覆盖、政策与上下文影响。因此本章先比较成对 prompt，再做 feature suppression，最后只提出范围有限的机制假说。我们关心的是哪些内部证据改变了输出倾向，不把某个 feature 的名字直接等同于系统级 truthfulness 或 safety。

\lecturefigure{slide-44-hallucinations-intro.jpg}{Parallel Processing: Hallucinations}{00:59:56}

\subsection{一个 always-on 的 IDK 倾向}

比较 Michael Jordan 与虚构的 Michael Batkin：前者激活 known-entity 与 basketball 证据，后者缺少可靠实体证据并生成不确定回答。追踪结果提出一个反直觉模型：不是“未知时才打开 IDK”，而是 IDK/refusal 倾向默认存在，已知识别 feature 会抑制它。

\lecturefigure{slide-45-always-on-idk-circuit.jpg}{已知实体通过 inhibition 压制默认的 IDK 路径}{00:59:59}

\lecturefigure{slide-46-knowing-unknowns.jpg}{已知与未知实体的输出对照}{01:00:17}

用一个极简 logit 模型表示：
\[
z_{\text{answer}}=b_{\text{answer}}+s_{\text{knowledge}},\qquad
z_{\text{IDK}}=b_{\text{IDK}}-\gamma s_{\text{known}},
\]
其中 $b$ 是默认偏置，$s_{\text{knowledge}}$ 是具体答案证据，$s_{\text{known}}$ 是“实体可识别”证据，$\gamma>0$ 表示抑制强度。真实网络更复杂；该式只表达“已知识别可压低 IDK”这一竞争关系。

\lecturefigure{slide-47-suppress-idk-hallucinates.jpg}{人为压制 IDK 回路后，未知实体被迫得到不稳定猜测}{01:02:22}

干预是关键：压低 IDK feature 后，模型对 Michael Batkin 给出 chess、tennis、hockey 等不同猜测。它说明该回路对拒答有因果作用，也说明“总让模型回答”会把不确定性转成看似具体的事实。

\lecturefigure{slide-48-natural-hallucination-hypothesis.jpg}{自然幻觉可能来自错误的 known-entity 信号压制了 IDK}{01:03:22}

图 48 把机制推广到真实名字与论文题目：某些表面熟悉 feature 可能错误地发出“我知道这个实体”的信号，从而解除 IDK 抑制，而具体答案证据又不足，于是模型用相关但错误的标题补全。这里的措辞必须是“some hallucinations might”，不能把所有错误都归结为一个 IDK 回路。

三张机制图之间存在严格的论证顺序。Michael Jordan 与 Michael Batkin 的对照先提出“已知/未知”区分；压制 IDK 的实验说明拒答方向对行为有因果作用；论文标题案例才把它作为自然错误的候选解释。第三步没有直接操纵现实世界中的全部知识，因此结论强度低于第二步。好的讲义必须把“观察到”“干预支持”“推测推广”分层书写，否则读者会把假说误读成已经定位所有幻觉根因。

\begin{warningbox}{幻觉不是单一机制，也不总是单个坏 token}
事实回忆、长文本一致性、工具结果误读、指令冲突和创造性续写都可能被称为 hallucination。这个案例解释的是一类实体问答中的拒答抑制；它不提供普遍 detector，也不定义所有场景中的幻觉边界。
\end{warningbox}

\subsection{Jailbreak 不是安全开关失灵，而是竞争路径接力}

直接问“怎样造炸弹”会立即激活 harmful-request 与 refusal 回路；acrostic prompt 先要求取首字母，模型在局部任务中输出 BOMB，随后 continuation 路径推动它继续解释。模型可能已经识别危险，却因为不同阶段的 feature 强度和路由不同而没有及时停止。

\lecturefigure{slide-49-jailbreaks-intro.jpg}{Parallel Processing: Jailbreaks}{01:03:47}

\lecturefigure{slide-50-jailbreak-comparison.jpg}{直接危险请求与首字母 jailbreak 的行为对照}{01:03:55}

\lecturefigure{slide-51-default-refusal-circuit.jpg}{默认路径：bomb 概念推动 harmful request 与 refusal}{01:04:28}

\lecturefigure{slide-52-jailbreak-keeps-going.jpg}{关键问题：模型已经写出 BOMB，为什么仍继续生成}{01:04:30}

读这组图时要比较 token 时序：prompt 阶段，acrostic/first-letter feature 占主导；生成 BOMB 后，危险概念与 refusal 开始增强；但语言 continuation 与“回答用户任务”路径已经形成惯性。安全系统若只在 prompt 开头分类一次，就可能错过生成过程中后出现的风险状态。

这也解释了“模型明明知道危险，为什么还继续”并非逻辑矛盾。神经网络可以同时含有危险识别与继续回答的证据，最终 token 由相对 logit 决定；识别到风险不保证 refusal 路径在每个时间点都占优。机制层的改进目标因此不只是增强一个 classifier，而是让风险 feature 在正确位置、正确时间持续影响解码，并在生成状态改变时允许重新路由或中止。

\begin{importantbox}{安全工程连接：把监控放进生成过程}
机制图提示三个实践方向：对中间状态持续监控，而非只做输入分类；允许模型在生成中途重新评估和停止；把“不知道/拒答”视为需要校准的竞争行为，而不是简单禁止或强制。
\end{importantbox}

\teachervoice{00:58:54--01:04:45，老师把 hallucination 与 jailbreak 都描述为并行路径：识别、回答、IDK、refusal 和 continuation 可以同时存在。课堂提示：机制结论是局部的，不能把一个漂亮回路图扩展成完整安全理论。}

\subsection{本章小结}

幻觉与 jailbreak 案例把“并行”从算术能力扩展到行为冲突。模型输出不是单模块决策，而是多个证据与抑制路径的合成。干预能证明某些路径有因果作用，却仍需区分案例机制、产品级风险和普遍安全结论。

